\section{Conclusion}

We have presented a controlled three-axis study of LLM-based and classical
survival prediction for all-cause mortality on a UK Biobank respiratory
disease subgroup ($n\,=\,124{,}158$).
By varying input representation (5 settings), encoder architecture (2
encoders), and feature fusion method (2 fusion strategies) under a uniform
CoxPH survival head, we isolate which structural encoding choices drive
prediction quality.

Three findings stand out.
First, \emph{temporal ordering is the primary discriminative structural
property}: decoupled age--event encoding (Setting 4) outperforms prose
serialisation of identical events by $+$0.008 C-index, regardless of
encoder or fusion.
Second, \emph{general-purpose LLMs implicitly encode numeric biomarker
structure}: a Qwen3-8B prose embedding without any raw numeric features
matches a fully supervised 39-feature CoxPH, with superior TD-AUC.
Third, \emph{autoregressive generation has a structural calibration
gap}: Delphi's near-term AUC collapse and high Brier score reveal that
pre-trained EHR LLMs without task-specific grounding are insufficient for
calibrated risk stratification.

Together, these findings confirm that \emph{how} temporal structure is
encoded matters more than which specific LLM encodes it, corroborated by
both the encoder ablation (Bio\_ClinicalBERT matches Qwen3-8B on Settings 3
and 4) and the fusion ablation (concatenation consistently outperforms
summation by $+$0.028--$+$0.039 C-index).
Future work will explore fine-tuning with calibration-aware objectives,
longer-context clinical encoders, and multimodal fusion of EHR embeddings
with imaging and genomic modalities.
